\documentclass[11pt]{article}
\usepackage[margin=1in]{geometry}
\usepackage{amsmath,amssymb,amsthm}
\usepackage{booktabs}
\usepackage[T1]{fontenc}
\usepackage[utf8]{inputenc}
\title{D5 note: T2 centered-core obstruction via four full collision families and one tail}
\author{}
\date{2026-03-21}

\newtheorem{theorem}{Theorem}
\newtheorem{proposition}[theorem]{Proposition}
\newtheorem{corollary}[theorem]{Corollary}
\newtheorem{remark}[theorem]{Remark}

\newcommand{\eX}{e_0}
\newcommand{\eQ}{e_1}
\newcommand{\eW}{e_2}
\newcommand{\eV}{e_3}
\newcommand{\eU}{e_4}

\begin{document}
\maketitle

\section*{Purpose}

This note closes the obstruction theorem
\[
T2:\qquad \sigma_{22}=[4,1,4]/[2,1,2]
\]
under the sharpened one-label tiny-repair background from d5\_226.
The claim is not merely that checked computation looks negative.
The point is that the surfaced 2026-03-21 bundle already exposes enough local
structure to write explicit collision families by hand.

The output is:
\begin{enumerate}
\item a row-level obstruction theorem for the base centered row $\sigma_{22}$,
      valid for every odd modulus $m\ge 5$;
\item an immediate corollary that the entire one-label tiny-repair scope of the
      endpoint artifact \texttt{032} (one-gate and one-bit repairs, with optional
      bundled cocycle-defect toggles) is already impossible for this row.
\end{enumerate}

Thus, if Background $T0$ is read exactly as the sharpened one-label repair scope,
then $T2$ is closed.

\section*{Coordinate conventions and the relevant background slices}

Write a color-$0$ vertex as
\[
x=(x_0,q,w,v,u)\in (\mathbb Z/m\mathbb Z)^5,
\]
and set
\[
\ell(x)=x_0+q+w+v+u \pmod m,
\qquad
s(x)=w+u \pmod m.
\]
Let $\eX,\eQ,\eW,\eV,\eU$ denote the standard coordinate increments in the
$(x_0,q,w,v,u)$ coordinates.

Let $F$ be the surfaced mixed background row coming from
\texttt{torus\_nd\_d5\_endpoint\_latin\_repair.py} via \texttt{\_mixed\_rule()}.
For the present argument we only need the following slices, all of which are read
from the surfaced tables $(0,2)$ on layer $2$ and $(0,3)/(3,0)$ on layer $3$:
\begin{enumerate}
\item every layer-$1$ source moves by $\eU$;
\item every layer-$2$ source with $q=m-1$ moves by $\eW$, while every layer-$2$
      source with $q\neq m-1$ moves by $\eX$;
\item every layer-$3$ source with $q=0$ and $s\neq 2$ moves by $\eX$.
\end{enumerate}
The background row $F$ is Latin.

\section*{The $\sigma_{22}$ row and its labelled source classes}

Fix the representative seed $(w_0,s_0)=(0,0)$ and set
\[
A:=\{x:\ \ell(x)=1,\ q=m-1,\ w=m-1,\ s(x)\neq 0\},
\]
\[
B:=\{x:\ \ell(x)=1,\ q=m-1,\ w=0,\ s(x)\neq 0\}.
\]
For
\[
\sigma_{22}=[4,1,4]/[2,1,2]
\]
the color-$0$ labelled source classes are
\[
L1=A,
\qquad
R1=B,
\qquad
L2=A+\eU,
\qquad
R2=B+\eW,
\qquad
L3=A+\eU+\eQ,
\qquad
R3=B+\eW+\eQ.
\]
Equivalently,
\[
\begin{array}{rcl}
L1&=&\{\ell=1,\ q=m-1,\ w=m-1,\ s\neq 0\},\\
R1&=&\{\ell=1,\ q=m-1,\ w=0,\ s\neq 0\},\\
L2&=&\{\ell=2,\ q=m-1,\ w=m-1,\ s\neq 1\},\\
R2&=&\{\ell=2,\ q=m-1,\ w=1,\ s\neq 1\},\\
L3&=&\{\ell=3,\ q=0,\ w=m-1,\ s\neq 1\},\\
R3&=&\{\ell=3,\ q=0,\ w=1,\ s\neq 1\}.
\end{array}
\]
Each class has cardinality $m(m-1)$.

Let $G$ denote the color-$0$ row defined by $\sigma_{22}$.
From the surfaced builder one reads:
\[
G=F\text{ on }L1,
\]
and on the remaining classes
\[
G(x)=x+\eW\text{ on }R1,
\qquad
G(x)=x+\eQ\text{ on }L2\cup R2,
\qquad
G(x)=x+\eU\text{ on }L3,
\qquad
G(x)=x+\eW\text{ on }R3.
\]
So the centered row differs from the background only on
\[
R1,\ L2,\ R2,\ L3,\ R3.
\]

\section*{Explicit collision families}

The key point is that one does not need the entire collision profile to obstruct
Latinness.  It is enough to write down a few disjoint overfull target families.

\begin{theorem}[explicit color-$0$ obstruction families for $T2$]
\label{thm:T2-families}
Let $m\ge 5$ be odd.
For the color-$0$ row $G$ of $\sigma_{22}=[4,1,4]/[2,1,2]$ at $(w_0,s_0)=(0,0)$,
there are five explicit background-vs-changed collision families:
\begin{center}
\begin{tabular}{@{}llll@{}}
\toprule
changed source $x$ & background mate $y$ & common target $t$ & count \\
\midrule
$x\in R1$ & $y=x+\eW-\eU$ & $t=x+\eW$ & $m(m-1)$ \\
$x\in L2$ & $y=x-\eX+\eQ$ & $t=x+\eQ$ & $m(m-1)$ \\
$x\in R2$ & $y=x-\eX+\eQ$ & $t=x+\eQ$ & $m(m-1)$ \\
$x\in R3$ & $y=x-\eX+\eW$ & $t=x+\eW$ & $m(m-1)$ \\
$x\in L3^0$ & $y=x-\eX+\eU$ & $t=x+\eU$ & $m$ \\
\bottomrule
\end{tabular}
\end{center}
where
\[
L3^0:=\{x\in L3:\ s(x)=0\}.
\]
For every row of the table,
\[
G(x)=F(y)=t,
\]
with $y$ a genuine background source.
The five target families are pairwise disjoint, so color $0$ already has at least
\[
4m(m-1)+m
\]
overfull targets.
Hence the full five-color row is non-Latin.
\end{theorem}

\begin{proof}
We treat the five families one by one.

\smallskip
\noindent\textbf{R1.}
Take $x\in R1$.
Then $x$ has layer $1$, $q=m-1$, $w=0$, and $s(x)\neq 0$.
Since $G$ uses the first right-step direction $2$, we have
\[
G(x)=x+\eW.
\]
Set $y=x+\eW-\eU$.
Then $y$ is still layer $1$, still has $q=m-1$, now has $w=1$, and therefore is
not in $L1$ or $R1$.
Every layer-$1$ background source moves by $\eU$, so
\[
F(y)=y+\eU=x+\eW=G(x).
\]
Different $x\in R1$ give different targets, hence this contributes $|R1|=m(m-1)$
distinct overfull targets.

\smallskip
\noindent\textbf{L2.}
Take $x\in L2$.
Then $x$ has layer $2$, $q=m-1$, $w=m-1$, and $s(x)\neq 1$.
Since $G$ uses the second left-step direction $1$, we have
\[
G(x)=x+\eQ.
\]
Set $y=x-\eX+\eQ$.
Then $y$ is still layer $2$, now has $q=0$, still has $w=m-1$, and therefore is
not in $L2$ or $R2$.
For layer $2$ with $q\neq m-1$, the background row uses $\eX$, so
\[
F(y)=y+\eX=x+\eQ=G(x).
\]
Again $x\mapsto x+\eQ$ is injective on $L2$, so this gives $|L2|=m(m-1)$
distinct overfull targets.

\smallskip
\noindent\textbf{R2.}
Take $x\in R2$.
Then $x$ has layer $2$, $q=m-1$, $w=1$, and $s(x)\neq 1$.
Again $G(x)=x+\eQ$.
Set $y=x-\eX+\eQ$.
Then $y$ is still layer $2$, now has $q=0$, still has $w=1$, and is therefore
background.
As above,
\[
F(y)=y+\eX=x+\eQ=G(x).
\]
This contributes $|R2|=m(m-1)$ distinct overfull targets.

\smallskip
\noindent\textbf{R3.}
Take $x\in R3$.
Then $x$ has layer $3$, $q=0$, $w=1$, and $s(x)\neq 1$.
Since $G$ uses the third right-step direction $2$, we have
\[
G(x)=x+\eW.
\]
Set $y=x-\eX+\eW$.
Then $y$ is still layer $3$, still has $q=0$, now has $w=2$, and
\[
s(y)=s(x)+1\neq 2
\]
because $s(x)\neq 1$.
Hence $y$ is a genuine background source on a layer-$3$, $q=0$, $s\neq 2$ slice,
so the background rule uses $\eX$ and
\[
F(y)=y+\eX=x+\eW=G(x).
\]
This contributes $|R3|=m(m-1)$ distinct overfull targets.

\smallskip
\noindent\textbf{The $L3$ tail.}
Take $x\in L3^0$, so $x\in L3$ and $s(x)=0$.
Then $x$ has layer $3$, $q=0$, $w=m-1$, and $G(x)=x+\eU$.
Set $y=x-\eX+\eU$.
Then $y$ is still layer $3$, still has $q=0$, still has $w=m-1$, and
\[
s(y)=1.
\]
But the changed class $L3$ is exactly the slice with $q=0$, $w=m-1$, and
$s\neq 1$.
So $y$ is background.
Since also $s(y)=1\neq 2$, the background row uses $\eX$ on $y$, hence
\[
F(y)=y+\eX=x+\eU=G(x).
\]
The set $L3^0$ has cardinality $m$: indeed $q=0$, $w=m-1$, and $s=0$ force $u=1$,
while the layer equation leaves one free parameter pair $(x_0,v)$ modulo one linear
constraint.
So this gives $m$ distinct overfull targets.

\smallskip
Finally, the target families are disjoint by their layer and $(q,w,s)$ data:
\[
\begin{array}{rcl}
R1&:& \ell=2,\ q=m-1,\ w=1,\\
L2&:& \ell=3,\ q=0,\ w=m-1,\\
R2&:& \ell=3,\ q=0,\ w=1,\\
R3&:& \ell=4,\ q=0,\ w=2,\\
L3^0&:& \ell=4,\ q=0,\ w=m-1,\ s=1.
\end{array}
\]
For $m\ge 5$ these slices are pairwise distinct.
Hence color $0$ has at least $4m(m-1)+m$ overfull targets, so the full row is
non-Latin.
\end{proof}

\begin{corollary}[row-level obstruction for the centered core]
For every odd $m\ge 5$, the centered row
\[
\sigma_{22}=[4,1,4]/[2,1,2]
\]
is non-Latin already at the base row.
\end{corollary}

\begin{proof}
Theorem~\ref{thm:T2-families} exhibits explicit overfull targets.
\end{proof}

\section*{One-label tiny repair is already impossible}

The endpoint artifact \texttt{032} allows two repair modes on top of the labelled
row:
\begin{enumerate}
\item a \emph{one-gate} repair, rerouting every source in one chosen label class;
\item a \emph{one-bit} repair, rerouting one chosen bit-slice inside one chosen
      label class.
\end{enumerate}
It also allows optional cocycle-defect toggles on two omitted layer-$3$ slices.
For $(w_0,s_0)=(0,0)$ these are
\[
O_L=\{\ell=3,\ q=m-1,\ w=0,\ s=2\},
\qquad
O_R=\{\ell=3,\ q=m-1,\ w=1,\ s=2\}.
\]
None of the explicit collision families above lies in these slices, and none of
their background mates lies there either.

\begin{proposition}[the one-label repair scope fails for $T2$]
\label{prop:T2-one-label-fails}
Let $m\ge 5$ be odd.
Take the centered row $\sigma_{22}=[4,1,4]/[2,1,2]$ and apply any one-gate or
one-bit repair from artifact \texttt{032}, with cocycle defect chosen from
\texttt{none}, \texttt{left}, \texttt{right}, or \texttt{both}.
The resulting row is still non-Latin.
\end{proposition}

\begin{proof}
A one-gate or one-bit repair changes sources inside at most one label class.
The cocycle-defect toggles change only the omitted slices $O_L$ and $O_R$, which
are disjoint from the explicit collision families and from their background mates.
Therefore the toggles do not touch the collisions from
Theorem~\ref{thm:T2-families}.

Among the explicit families, the four full families
\[
R1,\ L2,\ R2,\ R3
\]
live on four distinct labels.
A one-label repair can destroy at most one of them.
So after any one-label repair at least three full collision families survive.
In particular at least one explicit overfull target remains, hence the repaired row
is still non-Latin.
\end{proof}

\begin{remark}[scope and current status]
This proposition matches the exhaustive one-gate / one-bit negative scan from the
bundle, but the proof above does not use the full scan; it uses only the surfaced
row semantics plus the sharpened scope convention from d5\_226.
Therefore, once $T0$ is read as the one-label tiny-repair theorem, $T2$ is closed.
What remains outside $T0$--$T3$ is no longer a centered-strip obstruction problem,
but the later frontier package ($T4$, $G1$, $G2$).
\end{remark}

\begin{remark}[checked-range sanity certificate]
The companion checker
\texttt{d5\_228\_T2\_collision\_family\_checker.py} verifies the explicit family
formulae against the surfaced bundle on the checked range
\[
m=5,7,9,11,13.
\]
It confirms the expected counts
\[
|R1|=|L2|=|R2|=|R3|=m(m-1),
\qquad
|L3^0|=m,
\]
and reproduces the observed color-$0$ collision profile
\[
(B,R1),\ (B,L2),\ (B,R2),\ (B,R3)\text{ with multiplicity }m(m-1),
\qquad
(B,L3)\text{ with multiplicity }m.
\]
The proof above is symbolic; the checker is only a sanity certificate.
\end{remark}

\end{document}
